\section{StudyGraph: implementation and architecture}
\subsection{Problem, users, and MVP scope}
StudyGraph is for students revising \emph{LLMs for Developers}. Its value is course-specific assistance with inspectable sources and configurable teaching behaviour, rather than unrestricted general chat alone. The MVP provides three Streamlit pages: \textbf{Chat} for tasks and streamed answers with retrieved sources, \textbf{Buddies} for reusable teaching profiles, and \textbf{About} for model/index/storage information. Conversations and Buddy snapshots persist in SQLite. A Buddy can be course-linked (RAG) or unlinked (general chat). Explanations, summaries, practice questions, quizzes, comparisons, and answer review are included; unrestricted course management, student document upload, a learning-management system, and validated learning outcomes are outside scope.

\subsection{Architecture and user flow}
The application separates UI, services, domain models, orchestration, and external adapters. When a conversation starts, \sourcefile{ChatService} freezes the Buddy's Tone, Traits, Target, and course link into a snapshot, so later profile edits do not silently alter old conversations. For each message, history is loaded and a Task is selected. LangGraph routes course-linked turns through retrieval; general-chat turns skip it. Both then share prompt construction and Ollama generation before messages, prompt version, model, retrieval status, and actual source references are stored.

\begin{figure}[htbp]
\centering
\small
\fcolorbox{EKBlue}{LightBlue}{\parbox{0.96\linewidth}{\centering
Streamlit Chat $\rightarrow$ ChatService + frozen Buddy snapshot $\rightarrow$ LangGraph\\[2pt]
course-linked: current question $\rightarrow$ Ollama query embedding $\rightarrow$ Chroma top-$k$\\
general chat: retrieval skipped\\[2pt]
both: 4T prompt + history + any retrieved passages $\rightarrow$ local Ollama chat model\\
$\rightarrow$ streamed answer + source metadata $\rightarrow$ SQLite}}
\caption{Implemented StudyGraph turn flow; only the course-linked route retrieves.}
\label{fig:architecture}
\end{figure}

The defaults are \model{llama3.2:3b} for generation (temperature 0.2, context 4096) and \model{nomic-embed-text} for 768-dimensional embeddings. Both run locally through Ollama. Chroma and Streamlit run as separate local processes. This modular split and injected interfaces support tests without live services, although passing unit tests alone does not demonstrate correct generated answers.

\subsection{Course-material preparation and retrieval}
The source pipeline converts supplied PDFs/notebooks to cleaned Markdown; raw files are not indexed. The current corpus has 85 Markdown documents across eight teaching weeks, producing 1,660 chunks. Heading-first chunking keeps headings with content, targets 900 characters with 120-character overlap, drops chunks below 80 characters, and avoids splitting fenced code blocks. Deterministic IDs support idempotent upserts; collection metadata pins the embedding model, dimension, chunker version, and corpus digest to expose stale or incompatible indexes.

At query time the \emph{current user message} is embedded directly---there is no separate query-rewriting step. Chroma returns the nearest four passages by its configured vector distance. Their text, document/section metadata, and chunk IDs enter a delimited prompt block; only these retrieved sources can be shown as references. Retrieval distinguishes \texttt{ok}, \texttt{empty}, \texttt{failed}, and \texttt{skipped}. Empty or failed retrieval therefore cannot honestly be presented as grounded. A limitation remains: after a technical retrieval failure the graph still generates with empty context, so users must inspect status rather than treating every answer as RAG-grounded.

\subsection{Meaningful application of the 4T's}
The four dimensions are explicit fields rendered into \sourcefile{app/prompts/system\_v1.md}, not merely UI labels (Table~\ref{tab:4ts}). Task varies per interaction; the other three belong to the Buddy snapshot. Grounding rules additionally forbid invented references, treat retrieved text as data rather than instructions, require uncertainty when context is insufficient, and prohibit course citations in general chat.

\begin{table}[htbp]
\centering\small
\caption{Actual 4T implementation and purpose.}
\label{tab:4ts}
\begin{tabularx}{\linewidth}{@{}p{15mm}YY@{}}
\toprule
\textbf{Field} & \textbf{Where it appears} & \textbf{Why this choice}\\
\midrule
Traits & Buddy options such as patient, analytical, Socratic; inserted as ``Behave with these characteristics''. & Shapes teaching behaviour independently of topic.\\
Task & Per-message presets, e.g. explain, summarise, quiz, review; inserted under \texttt{\# Task}. & Makes the requested operation specific without creating a new Buddy per activity.\\
Tone & Buddy setting such as formal/academic or concise/direct; inserted under \texttt{\# Tone}. & Controls expression separately from the task.\\
Target & Learner level/goal, e.g. beginner or exam preparation; ``Address a learner described as\ldots''. & Adapts depth and framing to the intended student.\\
\bottomrule
\end{tabularx}
\end{table}
